\section{Results}\label{sec:results}

\FloatBarrier
\subsection{No module is retained on either specification}\label{sec:pooled}

Table~\ref{tab:rolling} gives the pooled rolling-origin scores and Fig.~\ref{fig:gaps} the persistence margins with their bootstrap intervals. On Specification A persistence has one-step RMSE 98.0\,kt and five-year RMSE 264.7\,kt; the modules range from 114.8 to 195.6\,kt at one year and from 288.6 to 488.2\,kt at five. On Specification B persistence has 84.4 and 300.0\,kt; the modules range from 119.5 to 205.7 and from 431.9 to 1058.9\,kt. No module comes within the 5\% band of persistence at either horizon on either specification; the closest approach is M1b on Specification A, 9.0\% adrift at five years and 17.1\% at one year. Hurdle (H2) therefore fails for every module, and (H1) adds nothing that could rescue it: the only comparator margin that excludes zero is M4 against M3 on Specification A at one year, and it runs the wrong way (M4 worse by 60.9\,kt, interval $[4.4, 134.4]$). The training-mean baseline is far worse than persistence everywhere (ratios 1.7--5.4) and is not discussed further.

\begin{table}[!htb]
\centering\small
\caption{Pooled rolling-origin scores (kt). RMSE with, in parentheses, the ratio to origin-matched persistence; MAE alongside. Specification A: coarse-regime catch, origins 1990--2014 ($n=25$, $h=1$) and 1990--2010 ($n=21$, $h=5$). Specification B: origins 1965--2023 ($n=59$) and 1965--2019 ($n=55$). A module is retained only if its ratio is below 0.95 at both horizons and it also beats its comparator (Table~\ref{tab:ladder}); none is.}
\label{tab:rolling}
\resizebox{\textwidth}{!}{%
\begin{tabular}{@{}lrrrrrrrr@{}}
\toprule
 & \multicolumn{4}{c}{Specification A} & \multicolumn{4}{c}{Specification B} \\
\cmidrule(lr){2-5}\cmidrule(lr){6-9}
 & \multicolumn{2}{c}{$h=1$} & \multicolumn{2}{c}{$h=5$} & \multicolumn{2}{c}{$h=1$} & \multicolumn{2}{c}{$h=5$} \\
Forecaster & RMSE & MAE & RMSE & MAE & RMSE & MAE & RMSE & MAE \\
\midrule
persistence & 98.0 (1.00) & 48.1 & 264.7 (1.00) & 137.2 & 84.4 (1.00) & 58.7 & 300.0 (1.00) & 217.3 \\
training mean & 423.6 (4.32) & 375.2 & 507.4 (1.92) & 447.8 & 458.1 (5.43) & 397.3 & 521.6 (1.74) & 456.8 \\
M1 autonomous & 120.5 (1.23) & 79.7 & 288.7 (1.09) & 158.2 & 119.5 (1.42) & 93.7 & 431.9 (1.44) & 335.7 \\
M1b depensatory & 114.8 (1.17) & 79.6 & 288.6 (1.09) & 158.2 & 151.6 (1.80) & 117.6 & 445.5 (1.49) & 344.1 \\
M2 stock-flow & 144.1 (1.47) & 60.6 & 398.2 (1.50) & 193.5 & 166.0 (1.97) & 106.2 & 1058.9 (3.53) & 758.7 \\
M3 AR residual & 134.6 (1.37) & 52.9 & 366.4 (1.38) & 183.2 & 126.7 (1.50) & 80.2 & 930.1 (3.10) & 650.8 \\
M4 lagged start & 195.6 (1.99) & 82.2 & 488.2 (1.84) & 223.6 & 205.7 (2.44) & 127.8 & 1030.7 (3.44) & 740.2 \\
\bottomrule
\end{tabular}}
\end{table}

\begin{figure}[!htb]
\centering
\includegraphics[width=\textwidth]{fig2_gaps}
\caption{Persistence margins: RMSE of each module minus RMSE of origin-matched persistence, as a percentage of the persistence RMSE, with 95\% paired moving-block bootstrap intervals (block $\max(h,3)$, 20\,000 replicates). The dashed line is the 5\% tie band; retention under (H2) requires the point estimate to lie below it. On Specification A every interval contains zero; on Specification B every interval lies above zero at the block length of the archived layer.}
\label{fig:gaps}
\end{figure}

The two specifications differ in what can be said about the margins. On Specification A every persistence margin has a bootstrap interval containing zero (M1 at one year: $+22.5\kt$, $[-13.4, +52.2]$, $p=0.19$; M3: $+36.6$, $[-16.8, +64.2]$, $p=0.93$), so the data do not show that the modules are worse than persistence, only that they are not demonstrably better, and the five-year M3 margin is the nearest to significance ($+101.6$, $[-0.8, +151.4]$, $p=0.08$). On Specification B every persistence margin excludes zero at the archived block length (M1 at one year: $+35.0$, $[+2.7, +70.8]$, $p=0.03$; M3: $+42.3$, $[+1.0, +92.5]$, $p=0.04$; all five-year margins $p\le 0.02$, M1 $[+33, +219]$). The M3 one-step margin on B is the only interval whose classification depends on the block length: it includes zero for blocks of four or more (SI Table S4); the others do not change. One property of the archived runner should be stated here because it affects the one-step margins: estimation pairs the transition $S_t\to S_{t+1}$ with the catch $C_t$, whereas the forecast path applies $C_{t+1}$ to the same step (Section~\ref{sec:scores}). Re-running the catch-supplied modules with the estimation convention, a sensitivity outside the archived protocol, reported in SI S14, lowers their one-step RMSE by 20--75\,kt (Specification A: M2 113.4, M3 105.2, M4 138.9\,kt; B: 146.4, 94.9, 130.3) and their five-year RMSE by up to 16\%; M3 then stands 7\% above persistence at one year on A and 12\% on B, still outside the band, and every five-year margin remains far outside it. The verdict does not depend on the convention; the size of the one-step deficit does. The MAE columns show that the one-step deficit is concentrated in a few large errors: on Specification A the M3 MAE is 10\% above persistence where its RMSE is 37\% above, and M2 is 26\% against 47\%. Where those errors sit is the subject of Section~\ref{sec:periods}.

\FloatBarrier
\subsection{Fixed windows: the collapse is missed by every module, the recovery is over-shot by the catch-supplied ones}\label{sec:fixed}

The fixed windows (Table~\ref{tab:fixed}, Fig.~\ref{fig:windows}) show the behaviour that the pooled scores summarise. From the 1983--1990 training window, in which SSB varied between 835 and 942\,kt, every module forecasts continuation. The autonomous modules place the stock at the fitted equilibrium, 888.6\,kt, and hold it there for five years; the catch-supplied modules, told that removals fall from 240 to 120 to 5\,kt, forecast growth to about 1050\,kt by 1993. The observed path falls from 862 to 9.7\,kt. Persistence (670\,kt RMSE) is the least bad because it at least does not grow; the autonomous modules (694) are close to it, and M2--M4 (819) are the worst because the catch cut, which in reality accompanied the collapse, is read by a stock-flow model as a reason for the stock to rise. On Specification B the pattern repeats from a 36-year training window (M1 817, M1b 906, M2 1898\,kt against persistence 633). Section~\ref{sec:identification} shows that this is not a fitting accident.

\begin{table}[!htb]
\centering\small
\caption{Fixed-window path RMSE (kt). Specification A, coarse-regime catch; the annual-landings pass and Specification B are given for the same windows. Persistence is frozen at the last training value; the training mean is the mean of the training window. Bold marks the lowest value in each column.}
\label{tab:fixed}
\begin{tabular}{@{}lrrrrrr@{}}
\toprule
 & \multicolumn{3}{c}{Specification A} & & \multicolumn{2}{c}{Specification B} \\
\cmidrule(lr){2-4}\cmidrule(lr){6-7}
 & collapse & recovery & recovery & & collapse & recovery stall \\
 & 1991--95 & 2008--15 & 2008--15 & & 1990--95 & 2013--24 \\
Forecaster & coarse & coarse & annual & & annual & annual \\
\midrule
persistence & \textbf{669.8} & 103.8 & 103.8 & & \textbf{632.7} & 261.7 \\
training mean & 687.9 & 144.4 & 144.4 & & 591.2 & 320.7 \\
M1 autonomous & 694.3 & 119.6 & 263.7 & & 817.3 & 254.3 \\
M1b depensatory & 693.6 & \textbf{89.8} & \textbf{78.1} & & 905.6 & \textbf{178.0} \\
M2 stock-flow & 819.4 & 119.6 & 302.8 & & 1897.7 & 268.8 \\
M3 AR residual & 819.4 & 219.7 & 609.1 & & 1767.8 & 267.7 \\
M4 lagged start & 819.1 & 214.1 & 586.3 & & 1871.9 & 268.8 \\
\bottomrule
\end{tabular}
\end{table}

\begin{figure}[!htb]
\centering
\includegraphics[width=\textwidth]{fig3_windows}
\caption{Fixed-window forecast paths on Specification A with the coarse-regime catch (the pass to which Table~\ref{tab:fixed} refers; the annual-landings pass is Fig.~S3). (a) Collapse window: every module forecasts continuation or growth from the 1990 state; M1b lies under M1. (b) Recovery window: M1 and M2 coincide (training-mean and realised catch are both close to 5\,kt), M1b is a flat line at its fitted $K$, and M3 and M4 over-predict most because the last training residual, $+28\kt$ in 2006--2007, is carried forward with $\phi = 0.95$.}
\label{fig:windows}
\end{figure}

The recovery window tells a different story. From 1995--2007, when SSB sat between 10 and 81\,kt, M1 and M2 coincide (training-mean catch and realised catch are both about 5\,kt) and over-predict the rebuild (120\,kt RMSE against 104 for persistence); M3 and M4 over-predict more (220 and 214), because the last training residual is a large positive increment ($+28\kt$ in 2006--2007) that the AR term with $\phi=0.95$ carries along the whole path. With annual landings the over-prediction is larger still (M3 609\,kt), since the fitted $K$ runs to its 5000-kt bound and production no longer saturates. The one module that beats persistence in this window, M1b (90 and 78\,kt), does so for a reason that undermines rather than supports it: its $\sA$ estimate descends to zero and its $K$ to 106--130\,kt, so it forecasts a flat line at a carrying capacity equal to the recent range. That is a persistence forecast produced by a production model whose parameters have lost their meaning, and the profile of its objective over $\sA$ is flat within 7\% (Fig.~S1). On Specification B the recovery-stall window shows the same ordering: M1b 178, M1 254, M2--M4 about 268\,kt, persistence 262.

\FloatBarrier
\subsection{Where the deficit arises: a period decomposition}\label{sec:periods}

The rule is applied to the pooled score, and the pooled score is what Table~\ref{tab:rolling} reports. Pooling, however, averages over regimes that the stock did not average over. Table~\ref{tab:periods} splits the one-step rolling errors into the periods visible in Fig.~\ref{fig:series}, and Fig.~\ref{fig:period} shows the per-origin errors behind it; this decomposition is post hoc and is reported as a description of the data, not as a second test. Three features emerge.

\begin{table}[!htb]
\centering\small
\caption{One-step rolling-origin RMSE (kt) by period of the forecast origin, with origin-matched persistence. Superscripts give the number of origins at which the module's forecast sits at the numerical floor ($<0.01\kt$), i.e.\ forecasts extinction; their absolute error is then the observed SSB. The post-1995 rows pool the later bands; bracketed values are paired bootstrap intervals for the module$-$persistence gap in kt (block 3, 20\,000 replicates; intervals are not shown for bands with $n<11$). Five-year values are in SI Table S5.}
\label{tab:periods}
\resizebox{\textwidth}{!}{%
\begin{tabular}{@{}llrrrrrrr@{}}
\toprule
Spec. & Origins & $n$ & persist & M1 & M1b & M2 & M3 & M4 \\
\midrule
A & 1990--1995 (collapse) & 6 & 193.6 & 208.8$^{3}$ & 186.4$^{4}$ & 292.0 & 273.2$^{3}$ & 397.3 \\
A & 1996--2012 (low biomass) & 17 & 19.6 & 64.8$^{12}$ & 69.7$^{13}$ & 17.2 & 14.7 & 19.4 \\
A & 2013--2014 (early rebuild) & 2 & 66.7 & 122.7 & 138.9 & 32.4 & 28.6 & 38.1 \\
A & 1996--2014 pooled & 19 & 28.5 & 73.1 & 79.8 & 19.4 & 16.7 & 22.1 \\
 & \quad gap vs.\ persistence & & & $+44.6$ & $+51.3$ & $-9.1$ & $-11.8$ & $-6.4$ \\
 & \quad 95\% interval & & & $[+29, +60]$ & $[+29, +69]$ & $[-17, +7]$ & $[-20, +3]$ & $[-14, +9]$ \\
\midrule
B & 1965--1990 (fished, pre-collapse) & 26 & 94.7 & 142.4$^{2}$ & 186.4$^{3}$ & 184.2 & 166.3 & 260.6 \\
B & 1991--1995 (collapse) & 5 & 167.0 & 126.8$^{3}$ & 66.6$^{3}$ & 299.8 & 154.8$^{3}$ & 321.6$^{1}$ \\
B & 1996--2012 (low biomass) & 17 & 20.7 & 75.0$^{17}$ & 75.0$^{17}$ & 26.3 & 19.7 & 24.0 \\
B & 2013--2023 (rebuild) & 11 & 61.0 & 112.3 & 174.5$^{1}$ & 160.5 & 96.4 & 136.3 \\
B & 1996--2023 pooled & 28 & 41.5 & 91.5 & 124.0 & 102.7 & 62.3 & 87.4 \\
 & \quad gap vs.\ persistence & & & $+50.0$ & $+82.5$ & $+61.2$ & $+20.9$ & $+46.0$ \\
 & \quad 95\% interval & & & $[+26, +79]$ & $[+53, +118]$ & $[-2, +93]$ & $[-11, +47]$ & $[-5, +73]$ \\
\bottomrule
\end{tabular}}
\end{table}

\begin{figure}[!htb]
\centering
\includegraphics[width=\textwidth]{fig4_period}
\caption{One-step absolute error by forecast origin (log scale, floored at 1\,kt for display) for persistence, M1, M2 and M3 on (top) Specification A and (bottom) Specification B. Shading marks the collapse band and the 1996--2012 low-biomass band. Open markers are forecasts at the numerical floor: at those origins M1's error is simply the observed SSB, which is why its curve traces the stock trajectory through 1996--2012 on both panels.}
\label{fig:period}
\end{figure}

First, the pooled deficit is carried by the collapse band. On Specification A the six origins 1990--1995 contribute 94\% of the pooled one-step sum of squared persistence errors and 63--99\% of each module's; the module-minus-persistence gap in that band runs from $-7\kt$ (M1b, a floor effect discussed below) to $+204\kt$ against a persistence RMSE of 194\,kt. Every forecaster is wrong by 130--590\,kt at the 1991 origin and by 140--710\,kt at the 1992 origin, because every forecaster, including persistence, extrapolates a state that is about to be removed. On Specification B the 26 fished pre-collapse origins add a second source of deficit: between 1965 and 1990 the catch-supplied modules are worse than persistence by 70--170\,kt, as a stock-flow model fitted to a series with 200--800\,kt removals mis-times the large swings of the 1960s and 1970s.

Second, the apparent collapse-band wins of M1, M1b and M3 on Specification B, and of M1b on A, are not forecasts of the collapse. At the 1993--1995 origins these modules forecast the floor: their training-mean catch (317--333\,kt on B, 177--208\,kt on A) exceeds the production of a stock of 10--100\,kt by an order of magnitude, so the map drives the state to $10^{-3}\kt$ in one step and the absolute error equals the observed value (10--31\,kt), which is small only because the stock was nearly gone. Only the 1991 and 1992 origins are genuine decline extrapolations, and there M1 beats persistence by 12 and 118\,kt on B (248 against 236, 273 against 155) while M2 and M4 are far worse. A floor forecast is a correct sign with no magnitude; it is credited by RMSE and should not be read as skill.

Third, the autonomous modules forecast extinction at every post-collapse origin, and the catch-supplied modules have one-step skill in the low-biomass regime on one specification only. Between 1996 and 2012 M1 and M1b sit at the floor at 17 of 17 origins on Specification B and at 12--13 of 17 on A (at every origin through 2007, and at all 15 five-year origins); their RMSE in that band, 65--75\,kt against 20\,kt for persistence, is the root mean square of the observed SSB. Training-mean catch is a plug, not an estimate, and after a collapse it is a plug an order of magnitude above sustainable removals; nothing in the autonomous branch can recover from that. In the same band M2, M3 and M4 forecast within 15--26\,kt, and on Specification A they beat persistence: pooling 1996--2014, M3 has one-step RMSE 16.7 against 28.5\,kt, a gap of $-11.8\kt$ ($-41\%$) with interval $[-20.5, +3.3]$, $p = 0.19$; M2 is $-32\%$ and M4 $-22\%$, and at five years M2 and M4 are $-30\%$ (SI Table S7). The intervals contain zero, the bands were chosen after the fact, and the post-1995 Specification A record is 19 one-step forecasts of a stock that moved slowly; but the point estimates are the right sign and the right order to be the residual signal that a catch-supplied production model should carry when the catch is small and the stock is far below $K$. On Specification B the same bands give the opposite answer: post-1995 every module is worse than persistence at one year (M3 by 50\%, M2 by 148\%), and at five years by factors of 2--6, because the 2013--2023 rebuild, which Specification A does not contain, is forecast by the stock-flow modules as a rise to several hundred kilotonnes above what occurred. The low-biomass skill does not transfer across the two series.

\FloatBarrier
\subsection{Variant passes}\label{sec:variants}

Four variant passes were added after the primary pass and declared in order (SI S1), each scored against persistence recomputed on its own origin set (Table~\ref{tab:variants}). Replacing the coarse catch regime by annual landings on Specification A leaves the ordering unchanged (M2 160.4, M3 153.6, M4 206.3\,kt at one year against 98.0 for persistence). Starting the stock-flow forecast from a survey-scaled index value instead of the assessed state gives 128.1 and 331.0\,kt. Allowing the intrinsic rate to differ before and after the 1991 capelin collapse \citep{Buren2014,DFO2024b}, the mechanism most often proposed for the post-collapse productivity regime, gives 153.8 and 334.0\,kt on A and 146.8 and 894.4 on B. Multiplying production by a capelin acoustic-index factor, which restricts the origins to the years with an acoustic estimate \citep{Murphy2025,NAFC2025}, gives 150.0 and 262.3\,kt on A against 96.6 and 193.0 for persistence on the reduced set, and 132.0 and 491.7 on B against 78.6 and 287.9. No variant comes within the band on either horizon.

\begin{table}[!htb]
\centering\small
\caption{Variant passes: rolling-origin RMSE (kt) of the variant module against persistence recomputed on the variant's origin set. $n$ is the number of origins at $h=1$ / $h=5$.}
\label{tab:variants}
\begin{tabular}{@{}llrrrr@{}}
\toprule
 & & \multicolumn{2}{c}{$h=1$} & \multicolumn{2}{c}{$h=5$} \\
\cmidrule(lr){3-4}\cmidrule(lr){5-6}
Variant & Spec.\ ($n$) & module & persist & module & persist \\
\midrule
Annual landings, M2 & A (25/21) & 160.4 & 98.0 & 393.8 & 264.7 \\
Annual landings, M3 & A (25/21) & 153.6 & 98.0 & 351.5 & 264.7 \\
Annual landings, M4 & A (25/21) & 206.3 & 98.0 & 486.4 & 264.7 \\
Survey-start M2 & A (25/21) & 128.1 & 98.0 & 331.0 & 264.7 \\
Two-regime $r$ (1991 break) & A (25/21) & 153.8 & 98.0 & 334.0 & 264.7 \\
Two-regime $r$ (1991 break) & B (59/55) & 146.8 & 84.4 & 894.4 & 300.0 \\
Capelin-index multiplier & A (24/20) & 150.0 & 96.6 & 262.3 & 193.0 \\
Capelin-index multiplier & B (36/32) & 132.0 & 78.6 & 491.7 & 287.9 \\
\bottomrule
\end{tabular}
\end{table}

\FloatBarrier
\subsection{Why the collapse cannot be forecast from inside the class}\label{sec:identification}

The fixed-window failure in Section~\ref{sec:fixed} has a structural explanation that does not depend on the fitted values. Write $F_C(S) = \Spos{S + g(S) - C}$ for the deterministic map at catch $C$.

\begin{proposition}\label{prop:repeller}
Let $g$ be continuous with $g(0)=0$, let $C>0$, and let $S_-(C) = \inf\{S>0 : g(S)\ge C\}$, with $S_-(C)=\infty$ if the set is empty.
\begin{enumerate}[label=(\roman*),leftmargin=2.2em,itemsep=1pt]
\item Every trajectory of $F_C$ started in $[0,S_-(C))$ is nonincreasing and remains in $[0,S_-(C))$; on any interval on which $g \ge C$, $F_C(S)\ge S$, so a trajectory cannot decline in one step from a state in that interval.
\item A one-step decline larger than the removals, $S_{t+1} < S_t - C_t$, requires $g(S_t) < 0$. On the compensatory branch this is impossible for $0\le S_t\le K$, whatever $(r,K)$; on the depensatory branch it requires $S_t < \sA$.
\end{enumerate}
\end{proposition}

The proof is in \ref{app:proof}. For the Schaefer branch with $4C\le rK$, $S_-(C) = \tfrac{K}{2}\bigl(1-\sqrt{1-4C/(rK)}\bigr)$ is the lower, unstable equilibrium and $S_+(C) = K - S_-(C)$ the upper, stable one, with $g \ge C$ between them; for $4C > rK$ production never reaches the catch and every trajectory declines to the floor. Three consequences follow for the ladder.

First, the lower equilibrium is generated by the catch: $S_-(C)\to 0$ as $C\to 0$. At the collapse-window fit ($r=1.935$, $K=1032.7\kt$; Fig.~\ref{fig:production}) it is 144\,kt under 240\,kt of removals, 66\,kt under 120\,kt and 2.6\,kt under 5\,kt. The threshold below which the autonomous map declines is a statement about sustained removals, not about the biology of small populations, and it is present in the compensatory branch with no depensation term at all. Second, the 1990 state, 862\,kt, lies between $S_- = 144$ and $S_+ = 888.6\kt$ for the fitted map under 240\,kt, so by (i) the autonomous forecast cannot decline at all; it rises to the attractor, which is what Fig.~\ref{fig:windows}a shows, and the attractor is locally stable ($F_C'(S_+) = -0.39$). Third, and independently of any fit, the 1991, 1992 and 1993 transitions are declines larger than the recorded catch: apparent net production $S_{t+1}-S_t+C_t$ is $-113$, $-161$ and $-66\kt$ under the coarse regime and $-181$, $-240$ and $-59\kt$ under annual landings (Fig.~\ref{fig:production}). By (ii) no parameter value of the compensatory branch, and no catch-supplied module built on it, can reproduce these steps; the depensatory branch can do so only with a threshold $\sA$ above 735\,kt, which no fit approaches. The collapse lies outside the model class at the recorded catch. The assessment model itself accommodates these years through under-reported catch and an elevated natural mortality \citep{Cadigan2016}; a forecaster that takes the catch record as given cannot, and a forecaster that is told the catch is falling to 5\,kt forecasts growth.

\begin{figure}[!htb]
\centering
\includegraphics[width=0.9\textwidth]{fig5_production}
\caption{Apparent net production $S_{t+1}-S_t+C_t$ against $S_t$ on Specification A, with the collapse-window fit ($r=1.935$, $K=1032.7$; maximum $rK/4 = 500\kt$ at $S=516$) and the two recovery-window fits. Under 240\,kt of removals the collapse-window map has its repeller at $S_- = 144\kt$ and its attractor at $S_+ = 889\kt$. The 1991--1993 transitions lie below every nonnegative production curve: the stock lost more biomass than the recorded catch removed. The recovery-window fits are nearly linear through the low-biomass points, with $K$ at its initialiser (500\,kt, coarse regime) or at its upper bound (5000\,kt, annual landings).}
\label{fig:production}
\end{figure}

The depensatory branch does not change this. In the fixed windows, which exclude the collapse transitions from training, M1b's $\sA$ estimate sits at its lower bound of zero, $r$ at its upper bound of 2 and $K$ at 1059\,kt (collapse window) or 106--130\,kt (recovery window), and the profile of the objective over $\sA$ changes by 7\% (coarse regime) and 3\% (annual landings) across its whole range (Fig.~S1). In rolling fits whose training window contains the collapse, the estimate jumps to 585--734\,kt on Specification A at every origin from 1992 to 2010, because the threshold has to sit above the collapse states for the branch to produce the negative production that (ii) requires, and it falls back to 5\,kt from 2011; on Specification B, where 36 pre-collapse years anchor the fit, it stays at zero throughout (SI S5). A depensation parameter that is zero, 700\,kt or 5\,kt depending on the window is not an estimate of an Allee threshold; it is the branch's only means of absorbing transitions the class cannot otherwise represent. Whether depensation has delayed the recovery \citep{Shelton1999,Hutchings1996,Rose2026} is therefore not a question this test can answer, and the usual difficulty of detecting depensation from assessment series \citep{Myers1995,Liermann1997} is compounded here by a collapse that the class attributes to it by default. The compensatory fits are weakly identified as well: $K$ rests at its 5000-kt bound at 23 of the 25 Specification A origins from 1992 onward and at 4800--5000\,kt on Specification B after 1993, so the post-collapse stock-flow forecaster is in effect linear, $S_{t+1}\approx(1+r)S_t - C$ with $r\approx 0.25$ (A) and $0.50$ (B), the familiar outcome of fitting a production model to a series with little contrast \citep{HilbornWalters1992,LudwigWalters1985,Polacheck1993}. That linear growth term is the whole of the low-biomass skill reported in Section~\ref{sec:periods}.

\FloatBarrier
\subsection{Operating characteristics of the rule}\label{sec:operating}

Fig.~\ref{fig:operating} reports the retention frequencies under the seven data-generating processes. When the truth is the autonomous Schaefer module at collapse-window parameters (D1), the rule retains M1 in 96.5\% and 98.5\% of replicates at the two noise levels: the empirical non-retention of M1 is not a power failure. At recovery-window parameters (D2) power is 71\% at the lower noise and 13\% at the higher; for the stock-flow truth (D3) it is 9--11\% and for the depensatory truth (D4) 0.5--1.5\%. Under the persistence-true null (D5) false retention is 1.5--3.0\%, within the bar. The two processes outside the class, added after the first round, fail the specificity bar badly: a drifting-productivity truth (D6) leads to retention of some module in 68--76\% of replicates and a smooth path observed with error (D7) in 93--98\%, most often M1. The rule is therefore specific against the absence of structure but not against the presence of structure of another kind; it detects that something beyond persistence is present, not that the retained module is that something.

\begin{figure}[!htb]
\centering
\includegraphics[width=\textwidth]{fig6_operating}
\caption{Retention frequencies of the rule under seven data-generating processes (200 replicates each, two process-noise levels; Specification A design). D1--D4: fraction of replicates in which the truth-matching module is retained (power; pre-declared bar 0.80). D5: fraction in which any module is retained under a persistence-true null (false retention; bar 0.10). D6--D7: the same under truths outside the ladder's class.}
\label{fig:operating}
\end{figure}

Table~\ref{tab:rules} evaluates alternative decision rules on the same replicates. The 5\% tie band carries most of the specificity (removing it costs 0.21 of specificity for 0.06 of power); the comparator hurdle costs 0.10 of power for 0.006 of specificity; a MASE-type scaled-error criterion alone \citep{HyndmanKoehler2006} is more powerful and far less specific; and an information criterion penalising parameter count dominates the adopted rule on both axes within the class. Two candidate diagnostics for predicting, from the data alone, whether the rule has power were examined and both fail (SI S8). The reading of the empirical verdict that these characteristics license is stated in Section~\ref{sec:discussion}.

\begin{table}[!htb]
\centering\small
\caption{Alternative decision rules evaluated on the archived simulation replicates. Mean power is averaged over D1--D4 and both noise levels; specificity is one minus the false-retention rate under D5, both noise levels. In-class truths only; under D6--D7 no rule examined is specific.}
\label{tab:rules}
\begin{tabular}{@{}lrr@{}}
\toprule
Decision rule & mean power & specificity \\
\midrule
Retention rule as adopted (H1--H3, 5\% band) & 0.376 & 0.978 \\
Without the comparator hurdle (H2, H3 only) & 0.476 & 0.972 \\
Persistence only, $h=1$, 5\% band & 0.562 & 0.955 \\
Persistence only, any margin & 0.542 & 0.765 \\
Full rule, no tie band & 0.436 & 0.772 \\
Full rule, 10\% band & 0.337 & 0.998 \\
MASE $<1$ against the naive benchmark & 0.651 & 0.675 \\
Information criterion, $n\log\mathrm{MSE} + 2k$ & 0.509 & 0.992 \\
\bottomrule
\end{tabular}
\end{table}

\FloatBarrier
\subsection{After 2015: which layer moved}\label{sec:verdict}

Specification A ends in 2015. The continuation of the same assessment family is published in the RAM Legacy database \citep[v4.66, stock COD2J3KL;][]{Ricard2012}, and xteNCAM reads the same years independently (Table~\ref{tab:outcome}). Three facts follow. The outcome window is an increase regime: every year 2016--2021 lies above its 2015 origin in both vintages, so it contains no event of the kind the collapse test probes, and persistence from the 2015 origin errs low throughout, with MAE 129.2\,kt and RMSE 133.3\,kt on the RAM vintage (107.5\,kt MAE from the Table A2 origin), the same order as its scored one-step RMSE. The two vintages agree on the outcome years to within 1--26\,kt (mean 16.8), while on the historical overlap they do not: the Table A2 vintage runs 24.7\,kt above the RAM vintage on average over 1983--2015 (mean absolute difference 26.8, maximum 131.2\,kt in 1984; Fig.~S2), so that the LRP defined as the 1983--1989 mean is 884.6\,kt on one vintage and 790.0\,kt on the other, a shift of 94.6\,kt in the reference point from a revision of the same model. And the two thresholds classify the same xteNCAM file differently in 29 of 71 years: under $\Blim = 276\kt$ the stock is below the limit in 26 years (1977--1979 and 1993--2015) and above it from 2016; under 884.6\,kt it has been below for 55 consecutive years since 1970 (SI S11). The 2023 framework revision, which moved the reported zone from critical to cautious without a change in the stock, is this reclassification executed by the assessment authority. Over the same period the observing system itself lapsed: partial fall-survey coverage in 2021 and 2022 left the leading indicator uncomputable, no typical assessment update was possible in 2022 or 2023, and the survey vessel was condemned (SI S12). None of this touches the RMSE comparison, which never uses a threshold; all of it bears on the verdict any forecast is scored against.

\begin{table}[!htb]
\centering\small
\caption{Outcome years on the Specification A horizon. Persistence error is $|S_t - S_{2015}|$ from the RAM-vintage 2015 origin of 277\,kt; the xteNCAM column is the revised model's reading of the same years.}
\label{tab:outcome}
\begin{tabular}{@{}lrrrr@{}}
\toprule
Year & RAM v4.66 (kt) & xteNCAM (kt) & $|$difference$|$ & persistence error (kt) \\
\midrule
2016 & 340 & 339 & 1 & 63 \\
2017 & 433 & 451 & 18 & 156 \\
2018 & 394 & 369 & 25 & 117 \\
2019 & 419 & 400 & 19 & 142 \\
2020 & 440 & 414 & 26 & 163 \\
2021 & 411 & 423 & 12 & 134 \\
\bottomrule
\end{tabular}
\end{table}
